\documentclass[12pt]{article}
\usepackage[T1]{fontenc}
\usepackage[utf8]{inputenc}
\usepackage[brazil]{babel}
\usepackage[margin=1in]{geometry}
\setlength{\parindent}{0pt}
\begin{document}
\section*{Tema 71}
Processo: PEDILEF 2006.71.50.010101-8/ RS\\
Situacao: Julgado\\
Ramo: DIREITO TRIBUTÁRIO\\
Questao: Saber se há incidência de imposto de renda sobre os benefícios e resgates da aposentadoria complementar sob a égide da Lei n. 7.713/88.\\
Tese: Os benefícios e resgates decorrentes das contribuições recolhidas sob o regime da Lei n. 7.713/88 (janeiro de 1989 a dezembro de 1995), com a incidência do imposto de renda no momento do recolhimento, não serão novamente tributados, sob pena de violação à regra proibitiva do "bis in idem". Vide recurso repetitivo do STJ: REsp 1001779/ DF.\\
Fonte: https://www.cjf.jus.br/phpdoc/virtus/
\end{document}
